% !PRINTED 176
\weberpage{176}

Wenn wir mit Hülfe der linearen Gleichung $b_1 = 0$ von
den fünf Grössen $\alpha$ die eine, etwa $\alpha_0$, durch die übrigen ausdrücken, so können wir die vier übrigen $\alpha_1, \alpha_2, \alpha_3, \alpha_4$ als homogene Coordinaten eines Raumpunktes betrachten. Die Gleichung
$b_2 = 0$ stellt dann eine Fläche zweiter Ordnung dar, $b_3 = 0$
eine Fläche dritter Ordnung, $b_4 = 0$ eine Fläche vierter Ordnung, $b_5 = 0$ eine Fläche fünfter Ordnung.

Wollte man also $b_2$, $b_3$, $b_4$ zugleich zu Null machen und
dadurch die gegebene Gleichung auf eine reine reduciren, so
müsste man einen der 24 Schnittpunkte von drei Flächen
zweiter, dritter, vierter Ordnung bestimmen, der von einer Gleichung $24^{\mathrm{ten}}$ Grades abhängt.

Statt nun so zu verfahren, sucht man auf der Fläche
$b_2 = 0$ eine gerade Linie zu bestimmen. Solcher gerader Linien
giebt es auf jeder Fläche zweiter Ordnung eine oder zwei
Schaaren (reell oder imaginär), und man kann eine dieser Geraden
\gesperrt{durch Quadratwurzeln} bestimmen, wie weiter unten ausgeführt werden soll. Die Ermittelung eines Schnittpunktes einer
dieser geraden Linien mit der Fläche $b_3 = 0$ führt dann auf
eine Gleichung dritten Grades, und wenn man noch durch Bestimmung eines gemeinsamen Factors der $\alpha$ den Coëfficienten
$b_4 = 1$ macht, so erhält die Gleichung für $y$ die Gestalt
\weq{6}{y^5 + y + b = 0,}
wodurch $y$ als Function \gesperrt{einer} Variabeln $b$ bestimmt ist.

Diese Gleichungsform wird gewöhnlich nach dem englischen
Mathematiker \gesperrt{Jerrard} genannt, der sie im Jahre 1834 bekannt
machte. Sie ist aber schon viel früher (1786) von \gesperrt{E. J. Bring}
in Lund publicirt worden\footnote{Vgl. F. Klein, Vorlesungen über das Ikosaëder, S. 143.}. Wir nennen sie daher die \gesperrt{Bring-Jerrard}'sche Form.

Ebensogut könnte man durch eine Gleichung vierten Grades
$b_4 = 0$ machen, und würde eine zweite Normalform der Gleichung
fünften Grades erhalten:
\weq{7}{y^5 + y^2 + b = 0,}
die die gleiche Berechtigung hat, wie die \gesperrt{Bring-Jerrard}'sche
Form.

Für die Lösung der Gleichung fünften Grades freilich wird
durch diese Betrachtungen nichts gewonnen. Ihr Nutzen besteht
%CONT>>
